\chapter{Summary and Conclusion}
\label{ch:conclusion}

\section{Summary of Contributions}
\label{sec:conclusion:summary}

This thesis developed methods for representing, generating, and evaluating
graphs. The work was motivated by a common difficulty across these tasks:
graphs are irregular, permutation-sensitive, and often constrained by
domain-specific validity requirements. The three contributions should be read
as operationally independent methods that address successive decisions in graph
generative modelling, rather than as a single fully automated end-to-end
system. They ask what graph information can be represented explicitly, how
graph-level information can guide generation, and how generated data can be
diagnosed after it is produced. The thesis made three main contributions.

\textbf{NSPPK for interpretable graph representation.}
The first contribution was the Neighborhood Subgraph Pairwise Path Kernel
(NSPPK), an explicit graph-to-vector representation for attributed graphs. For
each graph, NSPPK constructs a sparse feature vector by enumerating structural
patterns around pairs of anchor nodes. Each feature records the rooted
neighbourhoods around the anchors together with the shortest-path region that
connects them. In this way, the representation captures both local node-centred
structure and the broader context linking different parts of the graph.
Continuous node attributes are incorporated directly into these explicit
features, avoiding discretisation while retaining deterministic, sparse, and
interpretable feature construction.

Empirically, NSPPK showed that carefully designed explicit representations can
remain competitive with learned graph neural models, especially in low-data or
sample-efficient regimes. Its sparse feature map also supports computationally
transparent downstream learning and enables predictive behaviour to be related
back to structural graph patterns.

\textbf{GCDG for graph-conditioned generation.}
The second contribution was the Graph-Conditioned Diffusion Generator (GCDG), a
conditional graph generator that uses graph-level vectors to guide the
generation of node-level representations. The central contribution is the
graph-level-to-node-level conditional generation formulation: a fixed
graph-level descriptor conditions the generation of a padded node-level feature
matrix, from which the final graph is decoded. The diffusion process is defined
over continuous node-level features, while the denoising network predicts the
clean node representation and latent states used by downstream prediction
modules. These modules estimate node presence, node labels, degree-related
information, edge probabilities, and edge labels where available. A
constraint-aware decoder then combines the node-level and pairwise predictions
to construct the final discrete graph.

This staged design makes the generation process more inspectable than a single
black-box graph decoder. Intermediate quantities such as node-existence
probabilities, edge scores, degree predictions, and label predictions can be
examined before the final graph is selected. Experiments on synthetic
cycle--leg graphs and QM9 molecular graphs showed that conditioning quality and
graph validity are distinct issues: a model may produce plausible graphs while
failing to preserve the requested conditional structure. The chapter therefore
emphasised generation as conditional graph construction rather than
unconditional sampling alone.

\textbf{RankGen for statistically grounded evaluation.}
The third contribution was RankGen, a classifier-based framework for evaluating
generative models. RankGen treats evaluation as a collection of predictive
probes rather than as a single scalar similarity computation. The framework
defines four complementary metrics: Quality, Utility, Indistinguishability,
and Similarity, together with a composite score, Exchangeability. These metrics
diagnose different failure modes, including poor
task fidelity, lack of augmentation value, distributional separability, local
structural mismatch, and memorisation.

RankGen combines these probes with PAC-style finite-sample reliability
bounds and robust quantile-based aggregation. This separates the question of
whether a comparison is statistically reliable from the question of which
reliable model should be ranked higher. Experiments across controlled
perturbations, molecular graph
generators, image benchmarks, and open-domain text generation demonstrated that
the framework can expose behaviours that conventional metrics often collapse
or obscure.

\section{Key Insights and Lessons Learned}
\label{sec:conclusion:insights}

Several broader lessons emerge from these contributions.

First, explicit graph representations remain valuable. Although neural graph
models are highly expressive, deterministic feature maps can provide strong
performance, favourable sample efficiency, and clearer interpretability. NSPPK
shows that the limitations of classical kernels are not inherent to explicit
representations themselves; they can be addressed by enriching the structural
features and by treating continuous attributes directly.

Second, graph generation benefits from separating modelling roles. In GCDG, the
reverse diffusion process is applied to continuous node-level features, while
categorical node attributes, edge prediction, and final constraint satisfaction
are handled by dedicated components. This separation does not remove the
difficulty of graph generation, but it makes it easier to identify whether a
failure comes from the condition representation, the denoising network, the
prediction heads, or the decoder.

Third, conditional generation must be evaluated conditionally. It is not enough
for a generated graph to look plausible under aggregate statistics. If the
generator receives a graph-level condition, then the generated graph should be
tested against the information encoded in that condition. The synthetic
cycle--leg experiments illustrate this point clearly: different models can fail
in different conditional regimes even when their outputs remain graph-like.

Fourth, generative model evaluation is inherently multi-dimensional. A
generator may preserve predictive signal but be distributionally
distinguishable; it may generate realistic-looking samples that add little
utility; or it may perform well under global metrics while failing local
mixing or memorisation checks. RankGen addresses this by treating evaluation as
a structured statistical problem rather than a search for one universal score.

\section{Limitations}
\label{sec:conclusion:limitations}

The work in this thesis also has limitations.

For NSPPK, the use of explicit substructure enumeration provides
interpretability but can become expensive on dense or very large graphs. The
method is designed around bounded radii and distances, so its practical
scalability depends on parameter choices and graph density. The feature hashing
pipeline is efficient and deterministic, but hashed representations can still
merge distinct patterns when the hash space is limited. In addition, the
connector construction makes NSPPK more expressive than a purely local
neighbourhood-pair representation, but it also introduces an explicit trade-off:
larger connector radii and anchor distances can capture richer structure while
increasing feature extraction cost and sparsity.

For GCDG, the staged architecture introduces dependencies between components.
If the node representation does not preserve information needed for decoding,
the constrained decoder cannot recover it later. The decoder also introduces a
structural bias: because it enforces feasibility through an explicit graph
construction procedure, a weak or failed generator may still produce valid but
overly simple outputs, such as tree-like graphs, rather than visibly invalid
samples. For this reason, validity alone is not a sufficient measure of
generation quality; the generated graphs must also be evaluated for
condition preservation, structural statistics, labels, and distributional
agreement with the target graph family.

For RankGen, the framework depends on the quality of the chosen classifiers,
feature representations, and resampling protocol. Classifier-based evaluation
is powerful because it connects generation to predictive tasks, but it also
inherits the biases and limitations of the classifiers and embeddings used.
The PAC-style bounds provide conservative reliability checks, yet they do not
eliminate all sources of experimental variability. RankGen should therefore be
viewed as a diagnostic framework rather than as a replacement for all
domain-specific evaluation.

Across the thesis, a further limitation concerns the scope of the evaluation
formalism. RankGen is developed primarily around supervised, classifier-based
probes, where predictive accuracy and real-versus-generated discrimination
provide natural diagnostic signals. Extending the framework to regression
tasks, unsupervised learning, or settings without reliable labels would require
additional formal development, including appropriate probe definitions,
uncertainty bounds, and criteria for task-relevant utility.

Finally, the three methods are not used as a training loop. NSPPK is used as
the graph-level descriptor for GCDG, so representation and generation are
directly connected in Chapter~\ref{ch:generator}. RankGen could be used after
generation to compare GCDG with other generators, but in this thesis it is
developed and tested as a general evaluation framework in separate benchmark
settings. The thesis therefore supports a methodological link between
representation, generation, and evaluation, but does not claim that RankGen is
part of the GCDG training procedure.

\section{Future Research Directions}
\label{sec:conclusion:future}

Several directions follow naturally from this work.

\paragraph{Coupling explicit graph features with generators.}
One promising direction is to use NSPPK-like explicit graph features as an
interface between graphs and generators that do not operate directly on graph
objects. Many conditional generators require fixed-length vectors rather than
variable-size node and edge sets. In such settings, a precomputed NSPPK
descriptor could provide a graph-level representation that already encodes
neighbourhoods, attributes, and connector information. The generator would
therefore not have to learn this intermediate graph representation entirely
from data, but could focus on translating the supplied descriptor into
node-level and edge-level outputs.

\textbf{Scalable conditional graph generation.}
Future work on GCDG could improve scalability to larger and denser graphs by
developing more efficient decoding procedures, sparse pairwise prediction, and
hierarchical generation. Another direction is to integrate constraint handling
more closely with the denoising process, reducing the gap between neural
generation and final feasibility projection while retaining transparent
intermediate predictions.

\textbf{Richer condition preservation metrics.}
Conditional graph generation should be evaluated not only by asking whether
the generated graphs look realistic overall, but also by asking whether they
remain consistent with the graph-level information used as input to the
generator. Future work could therefore extend RankGen-style probes to measure
conditional dependence more directly: for example, whether quality and
diversity are maintained within each conditional regime, whether failures
concentrate in particular regions of the condition space, and whether changes
in the condition are reflected in the generated graphs without collapsing to a
single output pattern.

\textbf{Extending RankGen beyond classification.}
RankGen is formulated around classifier-based probes, where predictive accuracy
and real-versus-generated discrimination provide natural signals. Future work
could extend the same evaluation principle to regression, unsupervised
learning, or representation-learning tasks. This would require new probe
definitions and reliability checks, because the meaning of utility,
distributional alignment, and local similarity changes when class labels are
not available or when the downstream task is not classification.

\textbf{Molecular validation beyond graph statistics.}
For molecular graph generation, future evaluation should go beyond checking
whether the generated graphs match broad structural statistics. The generated
molecules should also be tested against chemistry-specific requirements such as
valency, allowed atom and bond types, synthesizability proxies, and the
property or scaffold constraints used as conditions. This would clarify whether
a conditional graph generator is producing chemically meaningful alternatives,
rather than only graphs that are valid under a generic graph decoder.

\section{Closing Remarks}
\label{sec:conclusion:closing}

The central message of this thesis is that graph generation cannot be judged by
the generator alone. Before a model can generate a graph, the graph or condition
has to be represented; after a graph is generated, its quality has to be tested.
These choices matter. The representation controls which structural information
is available to the generator, the generator controls how that information
becomes nodes, edges, and labels, and the evaluation protocol controls which
successes and failures are visible. The three parts therefore need to be
considered together, even when they are implemented as separate methods.

Three cross-cutting lessons follow. First, explicit graph features are useful
because they can often be inspected and related back to concrete graph
substructures. Their limitation is that they only capture the kinds of patterns
included in the feature design, so relevant structure outside that design may
be missed. Second, separating neural prediction from graph decoding can make
conditional generation easier to analyse, but it introduces decoder bias and
post-processing effects that validity metrics alone can hide. Third,
evaluation should not only ask whether generated graphs look reasonable on
average. It should also identify where they fail, how stable the comparison is,
and whether the outputs preserve the information used to condition generation.

The strongest claims supported by the thesis are therefore local and concrete:
NSPPK provides sparse attributed-graph features that can be inspected; GCDG
shows how graph-level descriptors can condition a staged graph generator; and
RankGen shows how predictive probes can expose failures that single scores
collapse. The broader lesson is methodological. Future graph generators should
be designed with their representations and evaluation protocols in view, so
that generated graphs are judged by whether they preserve the intended
information, respect relevant constraints, and remain useful under the tests
that matter for the target domain.
